% Part: first-order-logic
% Chapter: sequent-calculus
% Section: proving-things

\documentclass[../../../include/open-logic-section]{subfiles}

\begin{document}

\iftag{FOL}
      {\olfileid{fol}{seq}{pro}}
      {\olfileid{pl}{seq}{pro}}

\olsection{Voorbeelden van \usetoken{P}{derivation}}

\begin{ex}
Geef een $\Log{LK}$-!!{derivation} van de sequent $!A \land !B \Sequent !A$.

We beginnen door de gewenste eindsequent onderaan de !!{derivation}
te schrijven.
\begin{prooftree}
\AxiomC{}
\UnaryInf$!A\land !B \fCenter !A$
\end{prooftree}
Vervolgens moeten we bepalen welke afleidingsstap een onderste sequent
van deze vorm kan hebben. Dat zou een structurele regel kunnen zijn,
maar het is verstandig eerst naar een logische regel te zoeken. Het
enige logische connectief in de onderste sequent is $\land$. We zoeken
dus een $\land$-regel. Omdat het $\land$-teken in het antecedent
voorkomt, gaat het om de regel \LeftR{\land}.
\begin{prooftree}
\AxiomC{}
\RightLabel{\LeftR{\land}}
\UnaryInf$!A\land !B \fCenter !A$
\end{prooftree}
Er zijn twee mogelijkheden voor de bovenste sequent van deze
toepassing van \LeftR{\land}: dat kan $!A \Sequent !A$ of $!B
\Sequent !A$ zijn. De sequent $!A \Sequent !A$ is een beginsequent,
terwijl $!B \Sequent !A$ in het algemeen niet afleidbaar is. We vullen
daarom de eerste bovenste sequent in:
\begin{prooftree}
\Axiom$!A \fCenter !A$
\RightLabel{\LeftR{\land}}
\UnaryInf$!A\land !B \fCenter !A$
\end{prooftree}
Hiermee hebben we een correcte $\Log{LK}$-!!{derivation} van de
sequent $!A \land !B \Sequent !A$.
\end{ex}

\begin{ex}
Geef een $\Log{LK}$-!!{derivation} van de sequent $\lnot !A \lor !B
\Sequent !A \lif !B$.

Schrijf eerst de gewenste eindsequent onderaan de !!{derivation}.
\begin{prooftree}
\AxiomC{}
\UnaryInf$\lnot !A \lor !B \fCenter !A \lif !B$
\end{prooftree}
Om een logische regel te vinden die deze eindsequent kan opleveren,
bekijken we de logische connectieven erin: $\lnot$, $\lor$ en $\lif$.
Voorlopig zijn alleen $\lor$ en $\lif$ van belang, omdat zij de
!!{main operator}s zijn van !!{sentence}s in de eindsequent. Het
$\lnot$-teken staat binnen het bereik van een ander connectief en
komt daarom later aan bod. Voor de laatste afleidingsstap kunnen we dus
kiezen tussen \LeftR{\lor} en \RightR{\lif}. Beide regels zijn
mogelijk, maar we kiezen \RightR{\lif}, al was het maar omdat we het
opsplitsen in twee takken dan nog even uitstellen. Uit de vorm van een
toepassing van \RightR{\lif} met deze onderste sequent volgt dat de
afleiding er zo uit moet zien:
\begin{prooftree}
\AxiomC{}
\UnaryInf$ !A, \lnot !A \lor !B \fCenter !B $
\RightLabel{\RightR{\lif}} \UnaryInf$ \lnot !A \lor !B \fCenter !A \lif !B $
\end{prooftree}

Als we $\lnot !A \lor !B$ helemaal links in het antecedent plaatsen,
kunnen we de regel \LeftR{\lor} toepassen. Volgens het regelschema
moeten dan de volgende twee bovenste sequenten ontstaan:
\begin{prooftree}
\AxiomC{}
\UnaryInf$\lnot !A, !A \fCenter !B$
\AxiomC{}
\UnaryInf$!B, !A \fCenter !B$
\RightLabel{\LeftR{\lor}}
\BinaryInf$ \lnot !A \lor !B, !A \fCenter !B $
\RightLabel{\RightR{\Exchange}}
\UnaryInf$ !A, \lnot !A \lor !B \fCenter !B $
\RightLabel{\RightR{\lif}}
\UnaryInf$ \lnot !A \lor !B \fCenter !A \lif !B $
\end{prooftree}
Bedenk dat we van onder naar boven naar beginsequenten toewerken. We
zijn daar nu bijna. In de rechtertak zijn nog slechts één verzwakking
en één verwisseling nodig om een beginsequent te bereiken:
\begin{prooftree}
\AxiomC{}
\UnaryInf$\lnot !A, !A \fCenter !B$
\Axiom$!B \fCenter !B$
\RightLabel{\LeftR{\Weakening}}
\UnaryInf$!A, !B \fCenter !B$
\RightLabel{\LeftR{\Exchange}}
\UnaryInf$!B, !A \fCenter !B$
\RightLabel{\LeftR{\lor}}
\BinaryInf$\lnot !A \lor !B, !A \fCenter !B $
\RightLabel{\RightR{\Exchange}}
\UnaryInf$ !A, \lnot !A \lor !B \fCenter !B $
\RightLabel{\RightR{\lif}}
\UnaryInf$ \lnot !A \lor !B \fCenter !A \lif !B $
\end{prooftree}

In de linkertak is het enige logische connectief in !!{sentence} het
$\lnot$-teken in de !!{sentence}s van het antecedent. We hebben hier
dus een toepassing van de regel \LeftR{\lnot} nodig.
\begin{prooftree}
\AxiomC{}
\UnaryInf$ !A \fCenter !B, !A$
\RightLabel{\LeftR{\lnot}}
\UnaryInf$\lnot !A, !A \fCenter !B$
\Axiom$!B \fCenter !B$
\RightLabel{\LeftR{\Weakening}}
\UnaryInf$!A, !B \fCenter !B$
\RightLabel{\LeftR{\Exchange}}
\UnaryInf$!B, !A \fCenter !B$
\RightLabel{\LeftR{\lor}}
\BinaryInf$\lnot !A \lor !B, !A \fCenter !B $
\RightLabel{\RightR{\Exchange}}
\UnaryInf$ !A, \lnot !A \lor !B \fCenter !B $
\RightLabel{\RightR{\lif}}
\UnaryInf$ \lnot !A \lor !B \fCenter !A \lif !B $
\end{prooftree}
Net als in de rechtertak volstaan één verzwakking en één verwisseling
om ook de linkertak met een beginsequent af te sluiten.
\begin{prooftree}
\Axiom$!A \fCenter !A$
\RightLabel{\RightR{\Weakening}}
\UnaryInf$ !A \fCenter !A, !B$
\RightLabel{\RightR{\Exchange}}
\UnaryInf$ !A \fCenter !B, !A$
\RightLabel{\LeftR{\lnot}}
\UnaryInf$\lnot !A, !A \fCenter !B$
\Axiom$!B \fCenter !B$
\RightLabel{\LeftR{\Weakening}}
\UnaryInf$!A, !B \fCenter !B$
\RightLabel{\LeftR{\Exchange}}
\UnaryInf$!B, !A \fCenter !B$
\RightLabel{\LeftR{\lor}}
\BinaryInf$\lnot !A \lor !B, !A \fCenter !B $
\RightLabel{\RightR{\Exchange}}
\UnaryInf$ !A, \lnot !A \lor !B \fCenter !B $
\RightLabel{\RightR{\lif}}
\UnaryInf$ \lnot !A \lor !B \fCenter !A \lif !B $
\end{prooftree}
\end{ex}

\begin{ex}
Geef een $\Log{LK}$-!!{derivation} van de sequent $\lnot !A \lor \lnot !B
\Sequent \lnot (!A \land !B)$.

Met de voorgaande werkwijze schrijven we eerst de gewenste
eindsequent onderaan.
\begin{prooftree}
\AxiomC{}
\UnaryInf$ \lnot !A \lor \lnot !B \fCenter \lnot (!A \land !B) $
\end{prooftree}
De beschikbare hoofdconnectieven van de !!{sentence}s in de
eindsequent zijn $\lor$ en $\lnot$. Zowel \LeftR{\lor} als
\RightR{\lnot} kan hier worden toegepast. We beginnen met
\RightR{\lnot}, omdat we daarmee het opsplitsen in twee takken nog
even uitstellen:
\begin{prooftree}
\AxiomC{}
\UnaryInf$!A \land !B, \lnot !A \lor \lnot !B \fCenter $
\RightLabel{\RightR{\lnot}}
\UnaryInf$\lnot !A \lor \lnot !B \fCenter \lnot (!A \land !B)$
\end{prooftree}
Nu kunnen we kiezen tussen de regels \LeftR{\land} en
\LeftR{\lor}. Laten we eerst nagaan wat er gebeurt als we
\LeftR{\land} toepassen. Als bovenste sequent kunnen we dan $!A,
\lnot !A \lor !B \Sequent \quad$ of $!B, \lnot !A
\lor !B \Sequent \quad$ kiezen. De !!{derivation} is symmetrisch ten
opzichte van $!A$ en $!B$, dus kiezen we de eerste mogelijkheid:
\begin{prooftree}
\AxiomC{}
\UnaryInf$!A, \lnot !A \lor \lnot !B \fCenter $
\RightLabel{\LeftR{\land}}
\UnaryInf$!A \land !B, \lnot !A \lor \lnot !B \fCenter $
\RightLabel{\RightR{\lnot}}
\UnaryInf$\lnot !A \lor \lnot !B \fCenter \lnot (!A \land !B)$
\end{prooftree}
Als we de !!{derivation} verder invullen, lopen we echter vast:
\begin{prooftree}
\Axiom$!A \fCenter !A$
\RightLabel{\LeftR{\lnot}}
\UnaryInf$ \lnot !A, !A \fCenter$
\AxiomC{}
\RightLabel{?}
\UnaryInf$!A \fCenter !B$
\RightLabel{\LeftR{\lnot}}
\UnaryInf$ \lnot !B, !A \fCenter$
\RightLabel{\LeftR{\lor}}
\BinaryInf$\lnot !A \lor \lnot !B, !A \fCenter $
\RightLabel{\LeftR{\Exchange}}
\UnaryInf$!A, \lnot !A \lor \lnot !B \fCenter $
\RightLabel{\LeftR{\land}}
\UnaryInf$!A \land !B, \lnot !A \lor \lnot !B \fCenter $
\RightLabel{\RightR{\lnot}}
\UnaryInf$\lnot !A \lor \lnot !B \fCenter \lnot (!A \land !B)$
\end{prooftree}
De sequent bovenaan de rechtertak kan niet verder worden gereduceerd.
Ook kunnen structurele afleidingsstappen er geen beginsequent van
maken. Dit is dus niet de juiste route: de toepassing van
\LeftR{\land} hierboven was een vergissing. We keren terug naar de
voorgaande tussenstand en passen in plaats daarvan \LeftR{\lor} toe.
Dat geeft
\begin{prooftree}
\AxiomC{}
\UnaryInf$\lnot !A, !A \land !B \fCenter $

\AxiomC{}
\UnaryInf$\lnot !B, !A \land !B \fCenter $

\RightLabel{\LeftR{\lor}}
\BinaryInf$\lnot !A \lor \lnot !B, !A \land !B \fCenter $
\RightLabel{\LeftR{\Exchange}}
\UnaryInf$!A \land !B, \lnot !A \lor \lnot !B \fCenter $
\RightLabel{\RightR{\lnot}}
\UnaryInf$\lnot !A \lor \lnot !B \fCenter \lnot (!A \land !B)$
\end{prooftree}
Als we beide takken op de eerder gebruikte manier voltooien, krijgen
we
\begin{prooftree}
\Axiom$ !A \fCenter!A$
\RightLabel{\LeftR{\land}}
\UnaryInf$!A \land !B \fCenter !A$
\RightLabel{\LeftR{\lnot}}
\UnaryInf$\lnot !A, !A \land !B \fCenter $

\Axiom$ !B \fCenter !B$
\RightLabel{\LeftR{\land}}
\UnaryInf$!A \land !B \fCenter !B$
\RightLabel{\LeftR{\lnot}}
\UnaryInf$\lnot !B, !A \land !B \fCenter $

\RightLabel{\LeftR{\lor}}
\BinaryInf$\lnot !A \lor \lnot !B, !A \land !B \fCenter $
\RightLabel{\LeftR{\Exchange}}
\UnaryInf$!A \land !B, \lnot !A \lor \lnot !B \fCenter $
\RightLabel{\RightR{\lnot}}
\UnaryInf$\lnot !A \lor \lnot !B \fCenter \lnot (!A \land !B)$
\end{prooftree}
(We hadden in deze stappen de $\land$-regels ook lager dan de
$\lnot$-regels kunnen toepassen en toch een correcte !!{derivation}
verkregen.)
\end{ex}

\begin{ex}
Tot nu toe hebben we de contractieregel niet gebruikt, maar soms is
die wel nodig. Beschouw het volgende voorbeeld. We willen $\quad
\Sequent !A \lor \lnot !A$ bewijzen. Als we $\RightR{\lor}$
achterwaarts toepassen, krijgen we een van de volgende twee
!!{derivation}s:
\begin{prooftree}
\AxiomC{}
\UnaryInf$ \fCenter !A$
\RightLabel{\RightR{\lor}}
\UnaryInf$ \fCenter !A \lor \lnot !A$
\DisplayProof\qquad\bottomAlignProof
\AxiomC{}
\UnaryInf$!A \fCenter $
\RightLabel{\RightR{\lnot}}
\UnaryInf$ \fCenter \lnot !A$
\RightLabel{\RightR{\lor}}
\UnaryInf$ \fCenter !A \lor \lnot !A$
\end{prooftree}
Geen van beide loopt bovenaan uit op een beginsequent. Het beslissende
inzicht is dat we met de contractieregel twee voorkomens van
!!a{sentence} tot één kunnen samentrekken. Bij het zoeken naar een
bewijs, dus wanneer we van onder naar boven werken, kunnen we daarom
in de premisse een voorkomen van $!A \lor \lnot !A$ behouden,
bijvoorbeeld als volgt:
\begin{prooftree}
\AxiomC{}
\UnaryInf$ \fCenter !A \lor \lnot !A, !A$
\RightLabel{\RightR{\lor}}
\UnaryInf$ \fCenter !A \lor \lnot !A, !A \lor \lnot !A$
\RightLabel{\RightR{\Contraction}}
\UnaryInf$ \fCenter !A \lor \lnot !A$
\end{prooftree}
Nu kunnen we $\RightR{\lor}$ voor de tweede maal toepassen en ook
$\lnot !A$ verkrijgen. Zo komen we tot een volledige !!{derivation}.
\begin{prooftree}
\Axiom$!A \fCenter !A$
\RightLabel{\RightR{\lnot}}
\UnaryInf$\fCenter !A, \lnot !A$
\RightLabel{\RightR{\lor}}
\UnaryInf$\fCenter !A, !A \lor \lnot !A$
\RightLabel{\RightR{\Exchange}}
\UnaryInf$ \fCenter !A \lor \lnot !A, !A$
\RightLabel{\RightR{\lor}}
\UnaryInf$ \fCenter !A \lor \lnot !A, !A \lor \lnot !A$
\RightLabel{\RightR{\Contraction}}
\UnaryInf$ \fCenter !A \lor \lnot !A$
\end{prooftree}
\end{ex}

\begin{prob}
Geef !!{derivation}s van de volgende sequenten:
\begin{enumerate}
\item $!A \land (!B \land !C) \Sequent (!A \land !B) \land !C$.
\item $!A \lor (!B \lor !C) \Sequent (!A \lor !B) \lor !C$.
\item $!A \lif (!B \lif !C) \Sequent !B \lif (!A \lif !C)$.
\item $!A \Sequent \lnot\lnot !A$.
\end{enumerate}
\end{prob}

\begin{prob}
Geef !!{derivation}s van de volgende sequenten:
\begin{enumerate}
\item $(!A \lor !B) \lif !C \Sequent !A \lif !C$.
\item $(!A \lif !C) \land (!B \lif !C) \Sequent (!A \lor !B) \lif !C$.
\item $\Sequent \lnot(!A \land \lnot !A)$.
\item $!B \lif !A \Sequent \lnot !A \lif \lnot !B$.
\item $\Sequent (!A \lif \lnot !A) \lif \lnot !A$.
\item $\Sequent \lnot(!A \lif !B) \lif \lnot !B$.
\item $!A \lif !C \Sequent \lnot (!A \land \lnot !C)$.
\item $!A \land \lnot !C \Sequent \lnot (!A \lif !C)$.
\item $!A \lor !B, \lnot !B \Sequent !A$.
\item $\lnot !A \lor \lnot !B \Sequent \lnot(!A \land !B)$.
\item $\Sequent (\lnot !A \land \lnot !B) \lif\lnot(!A \lor !B)$.
\item $\Sequent \lnot(!A \lor !B) \lif (\lnot !A \land \lnot !B)$.
\end{enumerate}
\end{prob}

\begin{prob}
Geef !!{derivation}s van de volgende sequenten:
\begin{enumerate}
\item $\lnot(!A \lif !B) \Sequent !A$.
\item $\lnot(!A \land !B) \Sequent \lnot !A \lor \lnot !B$.
\item $!A \lif !B \Sequent \lnot !A \lor !B$.
\item $\Sequent \lnot \lnot !A \lif !A$.
\item $!A \lif !B, \lnot !A \lif !B \Sequent !B$.
\item $(!A \land !B) \lif !C \Sequent (!A \lif !C) \lor (!B \lif !C)$.
\item $(!A \lif !B) \lif !A \Sequent !A$.
\item $\Sequent (!A \lif !B) \lor (!B \lif !C)$.
\end{enumerate}
(Bij elk van deze opgaven is de regel $\RightR{\Contraction}$ nodig.)
\end{prob}
\end{document}
